\begin{section}{Metric Spaces}

Let $X$ be a set. A function $d: X \times X \to \R$ is called a \textbf{metric} on $X$ if for any $p, q \in X$, we have
\begin{enumerate}[label=\textbf{(\Alph*)}]
    \item $d(p, q) \geq 0$ and $d(p, q) = 0 \iff p = q$.
    \item $d(p, q) = d(q, p)$.
    \item $d(p, r) \leq d(p, q) + d(q, r)$ for any $r \in X$.
\end{enumerate}
\textbf{(C)} is called the \textbf{triangle inequality}. The pair $(X, d)$ is called a \textbf{metric space}.

\bigskip

\begin{theorem} \label{thm:StandardMetricOnRn}
    For some $n \in \N$, let $d: \R^{n} \times \R^{n} \to \R$ be a function defined by
    \[
    d(\bm{x}, \bm{y}) = \norm*{\bm{x} - \bm{y}} = \sqrt{\sum_{k = 1}^{n} (x_{k} - y_{k})^{2}}
    \]
    for all $\bm{x}, \bm{y} \in \R^{n}$. Then $d$ is a metric on $\R^{n}$.
    Especially, $d$ is called the \textbf{standard metric} on $\R^{n}$.
\end{theorem}

\begin{proof}
    Let $\bm{x}, \bm{y}, \bm{z} \in \R^{n}$ be given.
    \begin{itemize}
        \item $d(\bm{x}, \bm{y}) \geq 0$ is trivial. Also, we have
        \begin{align*}
            d(\bm{x}, \bm{y}) = 0 &\iff \norm*{\bm{x} - \bm{y}} = 0 \\
            &\iff x_{k} = y_{k} \text{ for all } k = 1, 2, \dots, n \\
            &\iff \bm{x} = \bm{y}.
        \end{align*}
        \item $d(\bm{x}, \bm{y}) = d(\bm{y}, \bm{x})$ is trivial.
        \item By \Cref{thm:PropertiesOfInnerProductNormDistance}, we have
        \begin{align*}
            d(\bm{x}, \bm{z}) &= \norm*{\bm{x} - \bm{z}} =
            \norm*{(\bm{x} - \bm{y}) + (\bm{y} - \bm{z})} \\
            &\leq \norm*{\bm{x} - \bm{y}} + \norm*{\bm{y} - \bm{z}}
            = d(\bm{x},\bm{y}) + d(\bm{y},\bm{z})
        \end{align*}
    \end{itemize}
    Therefore, $d$ is a metric on $\R^{n}$.
\end{proof}

\bigskip

\begin{definition} \label{def:KCell}
    Let $k \in \N$ be given. A set $E \subseteq \R^{k}$ defined by
    \[
    E \coloneqq \setbuilder*{\bm{x} = (x_{1}, x_{2}, \dots, x_{k})
    \in \R^{k}}{a_{j} \leq x_{j} \leq b_{j} \text{ for all } j = 1, 2, \dots, k}
    \]
    is called an \textbf{$k$-cell}, where $a_{j}, b_{j} \in \R$ with $a_{j} \leq b_{j}$ for all $j = 1, 2, \dots, k$.
\end{definition}

\bigskip

\begin{definition} \label{def:NeighborhoodOpenBall}
    Let $(X, d)$ be a metric space. For any $p \in X$ and $r \in \R^{+}$,
    we define the \textbf{neighborhood} of $p$ with radius $r$ as
    \[
    N_{r}(p) \coloneqq \setbuilder*{q \in X}{d(p, q) < r}.
    \]
    Also, we define the \textbf{punctured neighborhood} of $p$ with radius $r$ as
    \[
    Q_{r}(p) \coloneqq N_{r}(p) \setminus \set*{p} = \setbuilder*{q \in X}{0 < d(p, q) < r}.
    \]
    Especially, if $X = \R^{n}$ and $d$ is the standard metric on $\R^{n}$ for some $n \in \N$,
    then $N_{r}(p)$ is called the \textbf{open ball} of radius $r$ centered at $p$, which is denoted by $B_{r}(p)$.
\end{definition}

\bigskip

\begin{definition} \label{def:LimitIsolatedInteriorPoints}
    Let $(X, d)$ be a metric space and $E \subseteq X$. Let $p \in X$ be given.
    \begin{itemize}
        \item $p$ is called an \textbf{limit point} of $E$
        if for any $r \in \R^{+}$, we have $E \cap Q_{r}(p) \neq \emptyset$.
        \item $p$ is called a \textbf{isolated point} of $E$
        if there exists some $r \in \R^{+}$ such that $N_{r}(p) \cap E = \set*{p}$.
        \item $p$ is called an \textbf{interior point} of $E$
        if there exists some $r \in \R^{+}$ such that $N_{r}(p) \subseteq E$.
    \end{itemize}
\end{definition}

\bigskip

\begin{definition} \label{def:OpenClosedSets}
    Let $(X, d)$ be a metric space and $E \subseteq X$.
    \begin{itemize}
        \item $E$ is said to be \textbf{open} if any point of $E$ is an interior point of $E$.
        \item $E$ is said to be \textbf{closed} if any limit point of $E$ belongs to $E$.
    \end{itemize}
\end{definition}

\bigskip

\begin{definition} \label{def:BoundedPerfectDenseSets}
    Let $(X, d)$ be a metric space and $E \subseteq X$.
    \begin{itemize}
        \item $E$ is said to be \textbf{bounded} if there exists some $p \in X$ and
        $r \in \R^{+}$ such that $E \subseteq N_{r}(p)$.
        \item $E$ is said to be \textbf{perfect} if any point of $E$ is a limit point of $E$.
        \item $E$ is said to be \textbf{dense} if any point of $X$ is either a limit point of $E$ or a point of $E$.
    \end{itemize}
\end{definition}

\bigskip

\begin{theorem} \label{thm:NeighborhoodIsOpen}
    Let $(X, d)$ be a metric space. For any $p \in X$ and $r \in \R^{+}$, the neighborhood $N_{r}(p)$ is open.
\end{theorem}

\begin{proof}
    Let $q \in N_{r}(p)$ be given. Then we have $d(p, q) < r \implies \veps \coloneqq r - d(p, q) > 0$.
    Let $s \in N_{\veps}(q)$ be given. Then we have
    \[
    d(q, s) < \veps \implies d(p, s) \leq d(p, q) + d(q, s) < d(p, q) + \veps = r.
    \]
    Thus $s \in N_{r}(p)$, which implies that $N_{\veps}(q) \subseteq N_{r}(p)$. Since $q$ is arbtrary, $N_{r}(p)$ is open.
\end{proof}

\bigskip

\begin{theorem} \label{thm:LimitPointNeighborhoodIsInfinite}
    Let $(X, d)$ be a metric space and $E \subseteq X$. Let $p$ be a limit point of $E$.
    Then $E \cap N_{r}(p)$ is infinite for all $r \in \R^{+}$.
\end{theorem}

\begin{proof}
    Assume that there exists some $r \in \R^{+}$ such that $E \cap N_{r}(p)$ is finite. Let $S$ be a set defined by
    \[
    S \coloneqq (E \cap N_{r}(p)) \setminus \set*{p} = E \cap Q_{r}(p).
    \]
    Since $p$ is a limit point of $E$, we have $S \neq \emptyset$. Since $E \cap N_{r}(p)$ is finite,
    
    $S \subseteq E \cap N_{r}(p)$ is also finite by \Cref{thm:SubsetOfFiniteSetIsFinite}. Let $T$ be a set defined by
    \[
    T \coloneqq \setbuilder*{d(p, q)}{q \in S} \neq \emptyset.
    \]
    Since $p \notin S$, we have $0 \notin T$. Let $f: S \to T$ be a function defined by
    \[
    f(q) = d(p, q).
    \]
    Then $f$ is surjective. By \Cref{thm:SurjectionFromFiniteSetIsFinite}, $T$ is finite. Then, by
    \Cref{thm:FiniteOrderedSetHasExtremum}, $T$ has a minimum element, say $\veps = \min T > 0$. Then we have
    \begin{itemize}
        \item $N_{\veps}(p) \subset N_{r}(p)$.
        \item $E \cap Q_{\veps}(p) = \emptyset$.
    \end{itemize}
    This contradicts the assumption that $p$ is a limit point of $E$. Therefore, $E \cap N_{r}(p)$ is infinite.
\end{proof}

\bigskip

\begin{theorem} \label{thm:OpenSetComplementIsClosedSet}
    Let $(X, d)$ be a metric space and $E \subseteq X$. Then we have
    \[
    E \text{ is open} \iff E^{c} \text{ is closed.}
    \]
\end{theorem}

\begin{proof}
    \textbf{($\implies$)} Assume that $E$ is open. Let $p$ be a limit point of $E^{c}$.
    Then, for any $r \in \R^{+}$, we have
    \[
    E^{c} \cap Q_{r}(p) \neq \emptyset.
    \]
    If $p \in E$, then there exists some $\veps \in \R^{+}$ such that $N_{\veps}(p) \subseteq E$. Then we have
    \[
    E^{c} \cap Q_{\veps}(p) = \emptyset,
    \]
    which is a contradiction. Therefore, we have $p \in E^{c}$, which implies that $E^{c}$ is closed.
    \\[6mm]
    \textbf{($\impliedby$)} Assume that $E^{c}$ is closed. Let $p$ be a point of $E$.
    If $p$ is not an interior point of $E$, then for any $r \in \R^{+}$, we have
    \[
    E^{c} \cap N_{r}(p) \neq \emptyset.
    \]
    Since $p \in E \implies p \notin E^{c}$, we have
    \[
    E^{c} \cap Q_{r}(p) \neq \emptyset,
    \]
    which implies that $p$ is a limit point of $E^{c}$. Since $E^{c}$ is closed, we have $p \in E^{c}$,
    which is a contradiction. Therefore, $p$ is an interior point of $E$. Since $p$ is arbitrary, $E$ is open.
\end{proof}

\bigskip

\begin{corollary} \label{cor:ClosedSetComplementIsOpenSet}
    Let $(X, d)$ be a metric space and $E \subseteq X$. Then we have
    \[
    E \text{ is closed} \iff E^{c} \text{ is open.}
    \]
\end{corollary}